\documentclass[aspectratio=169]{beamer}
\usetheme{metropolis}
\usepackage{booktabs}
\usepackage{hyperref}

\title{Getting Started with AI Tools}
\subtitle{AI Bootcamp for UVM Economics --- Session 1}
\author{Emily Beam \and Erkmen Aslim}
\date{April 22, 2026}
\institute{Department of Economics, University of Vermont}

\begin{document}

\maketitle

% ============================================================
\section{Welcome}
% ============================================================

\begin{frame}{What is this?}
Two lunch sessions on AI tools for teaching and research.

\bigskip

\begin{itemize}
    \item Today: Getting started --- what AI can do, live demos
    \item April 27: Research pipelines, web scraping, automated grading
\end{itemize}

\bigskip
No installation required. Watch, ask questions, try things later.
\end{frame}

% ============================================================
\section{The AI Ladder}
% ============================================================

\begin{frame}{Where are you?}
\begin{table}
\begin{tabular}{cl}
\toprule
Level & Description \\
\midrule
0 & Haven't tried it \\
1 & Tried it once, wasn't impressed \\
2 & Use occasionally for simple tasks \\
3 & Regular user, integrated into workflow \\
4 & Power user, multi-step workflows \\
5 & Building systems with AI \\
\bottomrule
\end{tabular}
\end{table}

\bigskip
\textbf{Goal today:} Show you what Levels 3--4 look like.
\end{frame}

\begin{frame}{The tools landscape}
\begin{columns}
\begin{column}{0.5\textwidth}
\textbf{Chat interfaces}
\begin{itemize}
    \item ChatGPT, Claude.ai, Copilot
    \item Good for: quick questions, drafting
    \item Limit: copy-paste in/out
\end{itemize}
\end{column}
\begin{column}{0.5\textwidth}
\textbf{Terminal / Code tools}
\begin{itemize}
    \item Claude Code, GitHub Copilot
    \item Good for: file access, iteration, pipelines
    \item Limit: steeper learning curve
\end{itemize}
\end{column}
\end{columns}

\bigskip
UVM provides \href{https://go.uvm.edu/copilot}{Microsoft Copilot} free to everyone.
\end{frame}

% ============================================================
\section{Demo 1: Terminal + Claude Code (Emily)}
% ============================================================

% --- SHOW DURING DEMO ---
\begin{frame}{What you're seeing}
\begin{enumerate}
    \item \textbf{Ghostty} --- a terminal app (any terminal works)
    \item \textbf{Claude Code} --- AI assistant in the terminal
    \item Type a request $\rightarrow$ AI responds $\rightarrow$ iterate
\end{enumerate}

\bigskip
Key ideas:
\begin{itemize}
    \item AI can \textbf{read and write files} on your computer
    \item \textbf{Iterate}: ``Make it shorter,'' ``Add a section on X''
    \item Output is \textbf{Markdown} --- readable as plain text or rendered nicely
\end{itemize}
\end{frame}

\begin{frame}{Context and iteration}
\textbf{Context window}: The AI only ``sees'' what you show it. Be explicit.

\bigskip
\textbf{Iteration tips:}
\begin{itemize}
    \item First output is a draft, not a final product
    \item Refine: ``Change the tone,'' ``Add an example,'' ``This is wrong because\ldots''
    \item Save to a file so you can pick up later
\end{itemize}

\bigskip
\textbf{When it's wrong:} AI makes confident-sounding mistakes. Always review output, especially numbers and citations.
\end{frame}

% ============================================================
\section{Demo 2: Desktop App (Erkmen)}
% ============================================================

% --- Erkmen: add or replace slides here ---
\begin{frame}{Desktop AI tools}
% Placeholder --- Erkmen to fill in
\begin{itemize}
    \item ChatGPT / Claude desktop interface
    \item When to use chat vs.\ terminal
    \item Quick demo
\end{itemize}
\end{frame}

% ============================================================
\section{Demo 3: Teaching Application (Erkmen)}
% ============================================================

% --- Erkmen: add or replace slides here ---
\begin{frame}{AI for teaching tasks}
% Placeholder --- Erkmen to fill in
\begin{itemize}
    \item Creating assignment prompts
    \item Building grading rubrics
    \item Iterating on rubric criteria
\end{itemize}
\end{frame}

% ============================================================
\section{Wrap-up}
% ============================================================

\begin{frame}{Before next time (April 27)}
\begin{itemize}
    \item Try using an AI tool for \textbf{one real task} this week
    \item Think of a research or teaching task you'd like to automate
    \item Install guide: \url{https://eabeam.github.io/uvmecon-ai/prep.html}
\end{itemize}

\bigskip
\textbf{Next session:} Research pipelines, web scraping, automated grading
\end{frame}

\begin{frame}{Resources}
\begin{itemize}
    \item Bootcamp site: \url{https://eabeam.github.io/uvmecon-ai}
    \item Claude Code docs: \url{https://docs.anthropic.com/en/docs/claude-code}
    \item UVM Copilot (free): \url{https://go.uvm.edu/copilot}
    \item Anthropic prompt guide: \url{https://docs.anthropic.com/en/docs/build-with-claude/prompt-engineering/overview}
\end{itemize}
\end{frame}

\end{document}
